
\textbf{Corollaire 7.4.} \emph{Soit @@M0@@ un sch\'ema irr\'eductible de dimension @@M1@@, qui est projectif sur un corps alg\'ebriquement clos, et soit @@M2@@ un faisceau ample sur @@M3@@. Alors}~:
\begin{enumerate}[label=(\roman*)]
\item \emph{Pour qu'un faisceau inversible @@M4@@ sur @@M5@@ soit num\'eriquement \'equivalent \`a @@M6@@, (il faut et) il suffit que}
@@M7@@
\item \emph{Pour tout faisceau inversible @@M8@@ sur @@M9@@, qui n'est pas num\'eriquement proportionnel \`a @@M10@@, on a}
@@M11@@
\end{enumerate}

En effet, (ii) r\'esulte tout de suite de (i), d'apr\`es le raisonnement de 7.3. Pour \'etablir (i), soit @@M12@@ une courbe ferm\'ee, int\`egre, dans @@M13@@, on veut prouver que @@M14@@. Plongeons @@M15@@ dans @@M16@@ au moyen de @@M17@@ (@@M18@@). Soient @@M19@@ l'id\'eal homog\`ene de @@M20@@, et @@M21@@ un entier tel que @@M22@@ engendre @@M23@@ pour @@M24@@. Alors les hypersurfaces @@M25@@ de degr\'e @@M26@@ qui contiennent @@M27@@, n'ont pas de points communs en dehors de @@M28@@. Raisonnant par r\'ecurrence sur @@M29@@, on trouve \`a l'aide du th\'eor\`eme de Bertini @@M30@@ hypersurfaces @@M31@@ de degr\'e @@M32@@ qui contiennent @@M33@@, telles que @@M34@@ soit irr\'eductible. Comme @@M35@@, il suffit de prouver que @@M36@@ est num\'eriquement \'equivalent \`a @@M37@@ sur @@M38@@. Or par la formule de projection, on a que @@M39@@ et @@M40@@ ; donc l'assertion r\'esulte de 7.1.

\textbf{Corollaire 7.5.}
\emph{Soient @@M41@@ un corps alg\'ebriquement clos, @@M42@@ un @@M43@@-sch\'ema projectif, muni d'un faisceau ample @@M44@@, et @@M45@@ le sch\'ema des z\'eros d'une section @@M46@@ de @@M47@@. On suppose que pour toute composante irr\'eductible @@M48@@ de @@M49@@, 
